% ==============================================================================
%  CHAPTER 5 — CONCLUSION
% ==============================================================================
\section{Conclusion}
\label{sec:conclusion}

I set out to blur bystanders in first-person video while keeping the people the wearer is
with visible, using only the glasses' camera and without recognising anyone. This chapter
answers the three research questions from Section~\ref{sec:rq}, summarises the
contributions, and sets out what I would do next.

\subsection{Answers to the Research Questions}

\paragraph{RQ1: Can simple visual cues separate partners from bystanders without
identity recognition?}
Yes, with a clear bias towards privacy, and mostly through geometry. The size and position
of a face alone separated the two groups well enough that, on the video clips, the system
left only one of 33 bystanders visible, and 39 of the 40 faces it kept visible were genuine
partners. The ablation showed that head pose and mouth movement added almost nothing on top
of geometry in this dataset: the head-pose veto rarely had a case the geometry rule had not
already handled, and the speaking cue could only act on faces that already passed the size
test. The cues do not recognise everyone the wearer is with. Partners whose faces appear
small, such as friends across a table, fall below the size threshold, and some are
confirmed only after a delay. Neither failure needs identity information to fix.

\paragraph{RQ2: How well does selective anonymization protect bystanders and preserve
partners compared with blurring every face?}
Blanket blurring protects every bystander but keeps no partner visible. The live system
protected 96.97\% of bystanders while keeping 56.52\% of partners visible. For recorded
video, deciding once per person for the whole video, with the group rule and with the
pieces of each person's track joined by position and timing, raised partner preservation
to 85.51\% while still protecting 90.91\% of bystanders. Tested on held-out videos, letting the head-pose cue
lower the size threshold for faces turned towards the wearer raised the live system's
partner preservation to 78.26\% while blurring 90.91\% of bystanders, and keeping the
latch for as long as the tracker keeps a face raised it to 60.87\% at no cost to
bystanders. Selective anonymization therefore gives up a few
percentage points of bystander protection to recover most of what blanket blurring
destroys. These figures rest on labels from five annotators with almost perfect agreement,
held whichever annotator was taken as the reference and when whole videos were resampled,
but come from 21 videos, so they should be confirmed on more varied footage.

\paragraph{RQ3: How much computation does it need, and is it feasible on wearable-class
hardware?}
The classification itself costs almost nothing; the cost lies in face detection and,
for Phase~2b, in landmark extraction. With the detector exported to OpenVINO, the pipeline
runs in real time on a laptop CPU, at 31~FPS for Phase~2a and 22~FPS for Phase~2b, while
still finding 99.7\% of the faces the original detector found and making the same decision
on every labelled clip. On battery, the pipeline added about 5~W to the laptop's draw, or
about half a joule per frame. On the CPU of current
camera glasses alone it is projected at only 4 to 6~FPS, but with detection on the neural
processing unit these chips already contain, Phase~2a is projected at about 18~FPS, enough
for a live view, and Phase~2b at 8 to 12~FPS. Live operation on glasses is therefore
feasible, with landmark extraction as the remaining stage to move off the CPU. These device
figures are projections from published benchmarks and have not been measured on the
devices themselves.

\subsection{Contributions}

\begin{itemize}
  \item A selective anonymization pipeline that separates interaction partners from
        bystanders using face geometry, head pose and mouth movement from a single RGB
        camera, with no face recognition and no stored identity.
    \item An offline pass for recorded video that decides once per person (a retroactive
        latch, a group rule and track joining without face recognition) and keeps the people
        the wearer is with visible for their whole time on screen.
  \item An evaluation method for this task: blind labelling of short video clips with
        sound by several annotators, agreement measured with Cohen's and Fleiss' kappa,
        bystander protection and partner preservation reported separately with confidence
        intervals, every system decision checked against its rendered output, and a replay
        of the classifier on logged cues that measures each cue and threshold on the same
        labels.
  \item A breakdown of where the processing time goes, with projections to a
        Raspberry Pi~5 and to current camera-glasses hardware.
\end{itemize}

\subsection{Future Work}

\paragraph{Confirm the improvements on new data.}
Three improvements were tested after the main evaluation: the adaptive size threshold,
the persistent latch with an optional output delay, and the offline second pass. The
adaptive threshold was cross-validated across videos, but all three were developed with
knowledge of the same labelled clips. I would label a fresh set of clips from different
videos, with annotators who do not know the system, and measure them on it, then make the
best of them the default.

\paragraph{Move the landmark cues earlier in the rule.}
The ablation and the adaptive threshold together show that head pose is most useful when it
decides how small a face may be, not as a veto on faces that are already close. The next
version of the classifier should be designed around that, and the speaking cue should be
given more time to accumulate evidence on small faces.

\paragraph{Measure on real hardware.}
I would run the pipeline on a Raspberry Pi~5 with the detector exported to NCNN, and on
camera glasses with the detector on the NPU, measuring frame rate and power draw, and repeat
the labelled evaluation with detection on every second frame.

\paragraph{Broaden the dataset.}
The dataset lacks outdoor conversations and indoor crowds in which the wearer talks to no
one. Recording such scenes with consenting participants would also allow audio to be used
alongside the mouth cue, which would distinguish speaking from laughing or eating.
